%%%%%%%%%%%%%%%%%%%%%%%%%%%%%%%%%%%%%%%%%%%%%%%%%%%%%%%%%%%%
% 14.2 GREEK ALPHABET
% The Composition of Advanced Mathematics
%%%%%%%%%%%%%%%%%%%%%%%%%%%%%%%%%%%%%%%%%%%%%%%%%%%%%%%%%%%%

\section{Greek Alphabet}
\label{sec:greek-alphabet}

\prerequisite{
Basic familiarity with mathematical notation.
}

\learningobjectives{
By the end of this reference section, the reader should be able to:
\begin{itemize}
    \item identify all lowercase Greek letters;
    \item identify all uppercase Greek letters;
    \item distinguish Greek letters from similar Latin symbols;
    \item recognize common alternate Greek-letter forms;
    \item identify standard LaTeX commands for Greek letters;
    \item recognize common mathematical uses of Greek letters;
    \item distinguish notation conventions across mathematical subjects;
    \item interpret Greek letters used as variables, constants, parameters,
          angles, functions, measures, and operators;
    \item avoid common symbol-confusion errors;
    \item use Greek notation consistently in mathematical writing.
\end{itemize}
}

%%%%%%%%%%%%%%%%%%%%%%%%%%%%%%%%%%%%%%%%%%%%%%%%%%%%%%%%%%%%
\subsection{Why Greek Letters Are Used in Mathematics}
\label{subsec:why-greek-letters}
%%%%%%%%%%%%%%%%%%%%%%%%%%%%%%%%%%%%%%%%%%%%%%%%%%%%%%%%%%%%

Greek letters greatly expand the set of symbols available for mathematical
notation.

They are commonly used for:

\[
\boxed{
\text{angles},
}
\]

\[
\boxed{
\text{parameters},
}
\]

\[
\boxed{
\text{constants},
}
\]

\[
\boxed{
\text{eigenvalues},
}
\]

\[
\boxed{
\text{probability quantities},
}
\]

\[
\boxed{
\text{functions},
}
\]

and

\[
\boxed{
\text{operators}.
}
\]

The meaning of a Greek letter depends on context.

For example,

\[
\lambda
\]

may represent an eigenvalue, a Lagrange multiplier, a rate parameter, or a
wavelength.

%%%%%%%%%%%%%%%%%%%%%%%%%%%%%%%%%%%%%%%%%%%%%%%%%%%%%%%%%%%%
\subsection{Complete Greek Alphabet}
\label{subsec:complete-greek-alphabet}
%%%%%%%%%%%%%%%%%%%%%%%%%%%%%%%%%%%%%%%%%%%%%%%%%%%%%%%%%%%%

\begin{center}
\begin{tabular}{c|c|c|l}
\textbf{Name}
&
\textbf{Lowercase}
&
\textbf{Uppercase}
&
\textbf{Common LaTeX Commands}
\\ \hline

Alpha
&
\(\alpha\)
&
\(A\)
&
\verb|\alpha|
\\

Beta
&
\(\beta\)
&
\(B\)
&
\verb|\beta|
\\

Gamma
&
\(\gamma\)
&
\(\Gamma\)
&
\verb|\gamma|, \verb|\Gamma|
\\

Delta
&
\(\delta\)
&
\(\Delta\)
&
\verb|\delta|, \verb|\Delta|
\\

Epsilon
&
\(\epsilon,\varepsilon\)
&
\(E\)
&
\verb|\epsilon|, \verb|\varepsilon|
\\

Zeta
&
\(\zeta\)
&
\(Z\)
&
\verb|\zeta|
\\

Eta
&
\(\eta\)
&
\(H\)
&
\verb|\eta|
\\

Theta
&
\(\theta,\vartheta\)
&
\(\Theta\)
&
\verb|\theta|, \verb|\vartheta|, \verb|\Theta|
\\

Iota
&
\(\iota\)
&
\(I\)
&
\verb|\iota|
\\

Kappa
&
\(\kappa,\varkappa\)
&
\(K\)
&
\verb|\kappa|, \verb|\varkappa|
\\

Lambda
&
\(\lambda\)
&
\(\Lambda\)
&
\verb|\lambda|, \verb|\Lambda|
\\

Mu
&
\(\mu\)
&
\(M\)
&
\verb|\mu|
\\

Nu
&
\(\nu\)
&
\(N\)
&
\verb|\nu|
\\

Xi
&
\(\xi\)
&
\(\Xi\)
&
\verb|\xi|, \verb|\Xi|
\\

Omicron
&
\(o\)
&
\(O\)
&
Latin \(o,O\)
\\

Pi
&
\(\pi,\varpi\)
&
\(\Pi\)
&
\verb|\pi|, \verb|\varpi|, \verb|\Pi|
\\

Rho
&
\(\rho,\varrho\)
&
\(P\)
&
\verb|\rho|, \verb|\varrho|
\\

Sigma
&
\(\sigma,\varsigma\)
&
\(\Sigma\)
&
\verb|\sigma|, \verb|\varsigma|, \verb|\Sigma|
\\

Tau
&
\(\tau\)
&
\(T\)
&
\verb|\tau|
\\

Upsilon
&
\(\upsilon\)
&
\(\Upsilon\)
&
\verb|\upsilon|, \verb|\Upsilon|
\\

Phi
&
\(\phi,\varphi\)
&
\(\Phi\)
&
\verb|\phi|, \verb|\varphi|, \verb|\Phi|
\\

Chi
&
\(\chi\)
&
\(X\)
&
\verb|\chi|
\\

Psi
&
\(\psi\)
&
\(\Psi\)
&
\verb|\psi|, \verb|\Psi|
\\

Omega
&
\(\omega\)
&
\(\Omega\)
&
\verb|\omega|, \verb|\Omega|
\end{tabular}
\end{center}

%%%%%%%%%%%%%%%%%%%%%%%%%%%%%%%%%%%%%%%%%%%%%%%%%%%%%%%%%%%%
\subsection{Greek Letters With Latin-Looking Capitals}
\label{subsec:greek-latin-capitals}
%%%%%%%%%%%%%%%%%%%%%%%%%%%%%%%%%%%%%%%%%%%%%%%%%%%%%%%%%%%%

Several uppercase Greek letters are visually identical to Latin letters in
standard mathematical typesetting.

Examples include:

\[
\boxed{
A,\ B,\ E,\ Z,\ H,\ I,\ K,\ M,\ N,\ O,\ P,\ T,\ X.
}
\]

These correspond to uppercase forms of letters such as alpha, beta, epsilon,
zeta, eta, iota, kappa, mu, nu, omicron, rho, tau, and chi.

Because they do not provide visually distinct symbols, mathematics usually
uses the Latin character directly.

%%%%%%%%%%%%%%%%%%%%%%%%%%%%%%%%%%%%%%%%%%%%%%%%%%%%%%%%%%%%
\subsection{Alternate Lowercase Forms}
\label{subsec:alternate-greek-forms}
%%%%%%%%%%%%%%%%%%%%%%%%%%%%%%%%%%%%%%%%%%%%%%%%%%%%%%%%%%%%

Several Greek letters have commonly used alternate forms.

Important examples are

\[
\boxed{
\epsilon
\quad\text{and}\quad
\varepsilon,
}
\]

\[
\boxed{
\theta
\quad\text{and}\quad
\vartheta,
}
\]

\[
\boxed{
\pi
\quad\text{and}\quad
\varpi,
}
\]

\[
\boxed{
\rho
\quad\text{and}\quad
\varrho,
}
\]

\[
\boxed{
\sigma
\quad\text{and}\quad
\varsigma,
}
\]

and

\[
\boxed{
\phi
\quad\text{and}\quad
\varphi.
}
\]

These variants do not inherently have different mathematical meanings.

A text may assign them different meanings for convenience.

%%%%%%%%%%%%%%%%%%%%%%%%%%%%%%%%%%%%%%%%%%%%%%%%%%%%%%%%%%%%
\subsection{Alpha}
\label{subsec:greek-alpha}
%%%%%%%%%%%%%%%%%%%%%%%%%%%%%%%%%%%%%%%%%%%%%%%%%%%%%%%%%%%%

The lowercase letter alpha is

\[
\boxed{
\alpha.
}
\]

Common uses include:

\begin{itemize}
    \item angles;
    \item significance levels in statistics;
    \item parameters;
    \item coefficients;
    \item thermal expansion coefficients;
    \item regularization parameters;
    \item shape parameters.
\end{itemize}

For example, in hypothesis testing,

\[
\boxed{
\alpha
=
\Pr(\text{Type I error})
}
\]

under the standard testing framework.

%%%%%%%%%%%%%%%%%%%%%%%%%%%%%%%%%%%%%%%%%%%%%%%%%%%%%%%%%%%%
\subsection{Beta}
\label{subsec:greek-beta}
%%%%%%%%%%%%%%%%%%%%%%%%%%%%%%%%%%%%%%%%%%%%%%%%%%%%%%%%%%%%

The lowercase letter beta is

\[
\boxed{
\beta.
}
\]

Common uses include:

\begin{itemize}
    \item regression coefficients;
    \item distribution parameters;
    \item angles;
    \item discount parameters;
    \item epidemic transmission rates;
    \item financial market beta;
    \item special functions.
\end{itemize}

In regression,

\[
\boxed{
Y
=
X\beta+\varepsilon.
}
\]

%%%%%%%%%%%%%%%%%%%%%%%%%%%%%%%%%%%%%%%%%%%%%%%%%%%%%%%%%%%%
\subsection{Gamma}
\label{subsec:greek-gamma}
%%%%%%%%%%%%%%%%%%%%%%%%%%%%%%%%%%%%%%%%%%%%%%%%%%%%%%%%%%%%

The lowercase and uppercase forms are

\[
\boxed{
\gamma,
\qquad
\Gamma.
}
\]

Common uses include:

\begin{itemize}
    \item gamma distributions;
    \item Euler's gamma function;
    \item decay or recovery rates;
    \item paths or curves;
    \item Christoffel-symbol-related notation;
    \item option gamma in mathematical finance;
    \item Lorentz factors in physics.
\end{itemize}

The gamma function is

\[
\boxed{
\Gamma(s)
=
\int_0^\infty
x^{s-1}e^{-x}\,dx.
}
\]

For positive integers,

\[
\boxed{
\Gamma(n)
=
(n-1)!.
}
\]

%%%%%%%%%%%%%%%%%%%%%%%%%%%%%%%%%%%%%%%%%%%%%%%%%%%%%%%%%%%%
\subsection{Delta}
\label{subsec:greek-delta}
%%%%%%%%%%%%%%%%%%%%%%%%%%%%%%%%%%%%%%%%%%%%%%%%%%%%%%%%%%%%

The forms are

\[
\boxed{
\delta,
\qquad
\Delta.
}
\]

Lowercase delta often represents:

\begin{itemize}
    \item small changes;
    \item error tolerances;
    \item Kronecker delta;
    \item Dirac delta;
    \item decay parameters;
    \item depreciation rates.
\end{itemize}

Uppercase delta often represents:

\begin{itemize}
    \item finite changes;
    \item the Laplacian operator;
    \item discriminants;
    \item symmetric differences;
    \item option delta in finance.
\end{itemize}

For example,

\[
\boxed{
\Delta x
=
x_{\mathrm{final}}
-
x_{\mathrm{initial}}.
}
\]

%%%%%%%%%%%%%%%%%%%%%%%%%%%%%%%%%%%%%%%%%%%%%%%%%%%%%%%%%%%%
\subsection{Epsilon}
\label{subsec:greek-epsilon}
%%%%%%%%%%%%%%%%%%%%%%%%%%%%%%%%%%%%%%%%%%%%%%%%%%%%%%%%%%%%

Two common forms are

\[
\boxed{
\epsilon
}
\]

and

\[
\boxed{
\varepsilon.
}
\]

Epsilon commonly represents:

\begin{itemize}
    \item arbitrarily small positive quantities;
    \item approximation errors;
    \item residuals;
    \item strain in mechanics;
    \item tolerance values.
\end{itemize}

In the formal definition of a limit,

\[
\boxed{
\forall\varepsilon>0,
\quad
\exists\delta>0
}
\]

such that an appropriate closeness condition holds.

%%%%%%%%%%%%%%%%%%%%%%%%%%%%%%%%%%%%%%%%%%%%%%%%%%%%%%%%%%%%
\subsection{Zeta}
\label{subsec:greek-zeta}
%%%%%%%%%%%%%%%%%%%%%%%%%%%%%%%%%%%%%%%%%%%%%%%%%%%%%%%%%%%%

The lowercase letter is

\[
\boxed{
\zeta.
}
\]

Common uses include:

\begin{itemize}
    \item the Riemann zeta function;
    \item damping ratios;
    \item complex variables;
    \item random variables or parameters.
\end{itemize}

The Riemann zeta function is

\[
\boxed{
\zeta(s)
=
\sum_{n=1}^{\infty}
\frac{1}{n^s},
\qquad
\operatorname{Re}(s)>1.
}
\]

%%%%%%%%%%%%%%%%%%%%%%%%%%%%%%%%%%%%%%%%%%%%%%%%%%%%%%%%%%%%
\subsection{Eta}
\label{subsec:greek-eta}
%%%%%%%%%%%%%%%%%%%%%%%%%%%%%%%%%%%%%%%%%%%%%%%%%%%%%%%%%%%%

The lowercase letter eta is

\[
\boxed{
\eta.
}
\]

Common uses include:

\begin{itemize}
    \item efficiency;
    \item learning rates;
    \item viscosity parameters;
    \item random variables;
    \item coordinate parameters;
    \item Lagrange multipliers.
\end{itemize}

Its meaning is strongly context dependent.

%%%%%%%%%%%%%%%%%%%%%%%%%%%%%%%%%%%%%%%%%%%%%%%%%%%%%%%%%%%%
\subsection{Theta}
\label{subsec:greek-theta}
%%%%%%%%%%%%%%%%%%%%%%%%%%%%%%%%%%%%%%%%%%%%%%%%%%%%%%%%%%%%

Common forms are

\[
\boxed{
\theta,
\qquad
\vartheta,
\qquad
\Theta.
}
\]

Theta frequently represents an angle:

\[
\boxed{
\theta
=
\text{angle}.
}
\]

Other uses include:

\begin{itemize}
    \item unknown statistical parameters;
    \item model parameters;
    \item asymptotic notation;
    \item phase variables;
    \item polar coordinates.
\end{itemize}

In estimation,

\[
\boxed{
\hat{\theta}
}
\]

often denotes an estimator of parameter \(\theta\).

%%%%%%%%%%%%%%%%%%%%%%%%%%%%%%%%%%%%%%%%%%%%%%%%%%%%%%%%%%%%
\subsection{Iota}
\label{subsec:greek-iota}
%%%%%%%%%%%%%%%%%%%%%%%%%%%%%%%%%%%%%%%%%%%%%%%%%%%%%%%%%%%%

The lowercase letter is

\[
\boxed{
\iota.
}
\]

It is less common in elementary mathematics.

Typical uses include:

\begin{itemize}
    \item inclusion maps;
    \item embeddings;
    \item index-like parameters;
    \item interior products in differential geometry.
\end{itemize}

For an inclusion,

\[
\boxed{
\iota:A\hookrightarrow X.
}
\]

%%%%%%%%%%%%%%%%%%%%%%%%%%%%%%%%%%%%%%%%%%%%%%%%%%%%%%%%%%%%
\subsection{Kappa}
\label{subsec:greek-kappa}
%%%%%%%%%%%%%%%%%%%%%%%%%%%%%%%%%%%%%%%%%%%%%%%%%%%%%%%%%%%%

Common forms are

\[
\boxed{
\kappa,
\qquad
\varkappa.
}
\]

Kappa often represents:

\begin{itemize}
    \item curvature;
    \item condition numbers;
    \item thermal conductivity;
    \item diffusion coefficients;
    \item cardinal numbers;
    \item mean-reversion rates.
\end{itemize}

A matrix condition number may be written

\[
\boxed{
\kappa(A)
=
\|A\|
\|A^{-1}\|.
}
\]

%%%%%%%%%%%%%%%%%%%%%%%%%%%%%%%%%%%%%%%%%%%%%%%%%%%%%%%%%%%%
\subsection{Lambda}
\label{subsec:greek-lambda}
%%%%%%%%%%%%%%%%%%%%%%%%%%%%%%%%%%%%%%%%%%%%%%%%%%%%%%%%%%%%

The forms are

\[
\boxed{
\lambda,
\qquad
\Lambda.
}
\]

Lambda is one of the most frequently used Greek letters.

Common uses include:

\begin{itemize}
    \item eigenvalues;
    \item Lagrange multipliers;
    \item rates in Poisson and exponential distributions;
    \item wavelengths;
    \item regularization parameters;
    \item characteristic roots.
\end{itemize}

For eigenvalues,

\[
\boxed{
A\mathbf{v}
=
\lambda\mathbf{v}.
}
\]

%%%%%%%%%%%%%%%%%%%%%%%%%%%%%%%%%%%%%%%%%%%%%%%%%%%%%%%%%%%%
\subsection{Mu}
\label{subsec:greek-mu}
%%%%%%%%%%%%%%%%%%%%%%%%%%%%%%%%%%%%%%%%%%%%%%%%%%%%%%%%%%%%

The lowercase letter is

\[
\boxed{
\mu.
}
\]

Common uses include:

\begin{itemize}
    \item population means;
    \item measures;
    \item coefficients of friction;
    \item viscosity;
    \item growth or drift parameters;
    \item reduced mass;
    \item Lagrange multipliers.
\end{itemize}

For a population mean,

\[
\boxed{
\mu
=
\mathbb{E}[X].
}
\]

In measure theory,

\[
\boxed{
\mu(A)
}
\]

may denote the measure of set \(A\).

%%%%%%%%%%%%%%%%%%%%%%%%%%%%%%%%%%%%%%%%%%%%%%%%%%%%%%%%%%%%
\subsection{Nu}
\label{subsec:greek-nu}
%%%%%%%%%%%%%%%%%%%%%%%%%%%%%%%%%%%%%%%%%%%%%%%%%%%%%%%%%%%%

The lowercase letter is

\[
\boxed{
\nu.
}
\]

Common uses include:

\begin{itemize}
    \item frequency;
    \item degrees of freedom;
    \item Poisson ratio;
    \item kinematic viscosity;
    \item probability measures;
    \item parameters in special functions.
\end{itemize}

In fluid mechanics,

\[
\boxed{
\nu
=
\frac{\mu}{\rho}
}
\]

is kinematic viscosity.

%%%%%%%%%%%%%%%%%%%%%%%%%%%%%%%%%%%%%%%%%%%%%%%%%%%%%%%%%%%%
\subsection{Xi}
\label{subsec:greek-xi}
%%%%%%%%%%%%%%%%%%%%%%%%%%%%%%%%%%%%%%%%%%%%%%%%%%%%%%%%%%%%

The forms are

\[
\boxed{
\xi,
\qquad
\Xi.
}
\]

Xi is often used for:

\begin{itemize}
    \item random variables;
    \item coordinate variables;
    \item stochastic processes;
    \item auxiliary parameters;
    \item special functions.
\end{itemize}

%%%%%%%%%%%%%%%%%%%%%%%%%%%%%%%%%%%%%%%%%%%%%%%%%%%%%%%%%%%%
\subsection{Omicron}
\label{subsec:greek-omicron}
%%%%%%%%%%%%%%%%%%%%%%%%%%%%%%%%%%%%%%%%%%%%%%%%%%%%%%%%%%%%

Omicron is visually identical to the Latin letter \(o\):

\[
\boxed{
o.
}
\]

Its uppercase form is visually identical to

\[
\boxed{
O.
}
\]

Because of this, omicron is rarely used as a distinct mathematical symbol.

%%%%%%%%%%%%%%%%%%%%%%%%%%%%%%%%%%%%%%%%%%%%%%%%%%%%%%%%%%%%
\subsection{Pi}
\label{subsec:greek-pi}
%%%%%%%%%%%%%%%%%%%%%%%%%%%%%%%%%%%%%%%%%%%%%%%%%%%%%%%%%%%%

The forms are

\[
\boxed{
\pi,
\qquad
\varpi,
\qquad
\Pi.
}
\]

The most famous use is the circle constant

\[
\boxed{
\pi
\approx
3.141592653589793.
}
\]

It satisfies

\[
\boxed{
\pi
=
\frac{
\text{circumference}
}{
\text{diameter}
}.
}
\]

Uppercase pi is the product operator:

\[
\boxed{
\prod_{k=1}^{n}a_k.
}
\]

%%%%%%%%%%%%%%%%%%%%%%%%%%%%%%%%%%%%%%%%%%%%%%%%%%%%%%%%%%%%
\subsection{Rho}
\label{subsec:greek-rho}
%%%%%%%%%%%%%%%%%%%%%%%%%%%%%%%%%%%%%%%%%%%%%%%%%%%%%%%%%%%%

Common forms are

\[
\boxed{
\rho,
\qquad
\varrho.
}
\]

Rho frequently represents:

\begin{itemize}
    \item density;
    \item correlation;
    \item radial coordinates;
    \item spectral radius;
    \item discount rates;
    \item charge density.
\end{itemize}

For correlation,

\[
\boxed{
\rho_{X,Y}
=
\frac{
\operatorname{Cov}(X,Y)
}{
\sigma_X\sigma_Y
}.
}
\]

In continuum mechanics,

\[
\boxed{
\rho
=
\text{mass density}.
}
\]

%%%%%%%%%%%%%%%%%%%%%%%%%%%%%%%%%%%%%%%%%%%%%%%%%%%%%%%%%%%%
\subsection{Sigma}
\label{subsec:greek-sigma}
%%%%%%%%%%%%%%%%%%%%%%%%%%%%%%%%%%%%%%%%%%%%%%%%%%%%%%%%%%%%

The forms are

\[
\boxed{
\sigma,
\qquad
\varsigma,
\qquad
\Sigma.
}
\]

Lowercase sigma often represents:

\begin{itemize}
    \item standard deviation;
    \item stress;
    \item volatility;
    \item scale parameters;
    \item permutations.
\end{itemize}

Uppercase sigma represents summation:

\[
\boxed{
\sum_{i=1}^{n}a_i.
}
\]

It may also represent a covariance matrix:

\[
\boxed{
\Sigma
=
\operatorname{Cov}(\mathbf{X}).
}
\]

%%%%%%%%%%%%%%%%%%%%%%%%%%%%%%%%%%%%%%%%%%%%%%%%%%%%%%%%%%%%
\subsection{Tau}
\label{subsec:greek-tau}
%%%%%%%%%%%%%%%%%%%%%%%%%%%%%%%%%%%%%%%%%%%%%%%%%%%%%%%%%%%%

The lowercase letter is

\[
\boxed{
\tau.
}
\]

Common uses include:

\begin{itemize}
    \item time constants;
    \item stopping times;
    \item shear stress;
    \item torque;
    \item topology notation;
    \item optical depth;
    \item delay.
\end{itemize}

A delay differential equation may contain

\[
\boxed{
x(t-\tau).
}
\]

%%%%%%%%%%%%%%%%%%%%%%%%%%%%%%%%%%%%%%%%%%%%%%%%%%%%%%%%%%%%
\subsection{Upsilon}
\label{subsec:greek-upsilon}
%%%%%%%%%%%%%%%%%%%%%%%%%%%%%%%%%%%%%%%%%%%%%%%%%%%%%%%%%%%%

The forms are

\[
\boxed{
\upsilon,
\qquad
\Upsilon.
}
\]

Upsilon appears less frequently in general mathematics but is used in:

\begin{itemize}
    \item particle physics;
    \item special functions;
    \item auxiliary variables;
    \item advanced mathematical notation.
\end{itemize}

%%%%%%%%%%%%%%%%%%%%%%%%%%%%%%%%%%%%%%%%%%%%%%%%%%%%%%%%%%%%
\subsection{Phi}
\label{subsec:greek-phi}
%%%%%%%%%%%%%%%%%%%%%%%%%%%%%%%%%%%%%%%%%%%%%%%%%%%%%%%%%%%%

Common forms are

\[
\boxed{
\phi,
\qquad
\varphi,
\qquad
\Phi.
}
\]

Phi frequently denotes:

\begin{itemize}
    \item angles;
    \item functions;
    \item potentials;
    \item the golden ratio;
    \item the standard normal CDF;
    \item Euler's totient function;
    \item mappings and homomorphisms.
\end{itemize}

The golden ratio is often written

\[
\boxed{
\phi
=
\frac{1+\sqrt5}{2}.
}
\]

In probability,

\[
\boxed{
\Phi(x)
=
\Pr(Z\leq x),
\qquad
Z\sim N(0,1).
}
\]

%%%%%%%%%%%%%%%%%%%%%%%%%%%%%%%%%%%%%%%%%%%%%%%%%%%%%%%%%%%%
\subsection{Chi}
\label{subsec:greek-chi}
%%%%%%%%%%%%%%%%%%%%%%%%%%%%%%%%%%%%%%%%%%%%%%%%%%%%%%%%%%%%

The lowercase letter is

\[
\boxed{
\chi.
}
\]

Common uses include:

\begin{itemize}
    \item chi-square distributions;
    \item characteristic functions in some contexts;
    \item Euler characteristics;
    \item indicator-like functions;
    \item susceptibility in physics.
\end{itemize}

The chi-square statistic is commonly written

\[
\boxed{
\chi^2.
}
\]

%%%%%%%%%%%%%%%%%%%%%%%%%%%%%%%%%%%%%%%%%%%%%%%%%%%%%%%%%%%%
\subsection{Psi}
\label{subsec:greek-psi}
%%%%%%%%%%%%%%%%%%%%%%%%%%%%%%%%%%%%%%%%%%%%%%%%%%%%%%%%%%%%

The forms are

\[
\boxed{
\psi,
\qquad
\Psi.
}
\]

Psi commonly represents:

\begin{itemize}
    \item wave functions;
    \item stream functions;
    \item special functions;
    \item mappings;
    \item state functions.
\end{itemize}

In two-dimensional fluid dynamics, a stream function may be written

\[
\boxed{
\psi(x,y).
}
\]

%%%%%%%%%%%%%%%%%%%%%%%%%%%%%%%%%%%%%%%%%%%%%%%%%%%%%%%%%%%%
\subsection{Omega}
\label{subsec:greek-omega}
%%%%%%%%%%%%%%%%%%%%%%%%%%%%%%%%%%%%%%%%%%%%%%%%%%%%%%%%%%%%

The forms are

\[
\boxed{
\omega,
\qquad
\Omega.
}
\]

Lowercase omega commonly represents:

\begin{itemize}
    \item angular frequency;
    \item angular velocity;
    \item differential forms;
    \item vorticity;
    \item probability outcomes;
    \item natural frequencies.
\end{itemize}

Uppercase omega commonly represents:

\begin{itemize}
    \item sample spaces;
    \item domains;
    \item solid angles;
    \item asymptotic lower bounds.
\end{itemize}

In probability,

\[
\boxed{
\Omega
=
\text{sample space}.
}
\]

%%%%%%%%%%%%%%%%%%%%%%%%%%%%%%%%%%%%%%%%%%%%%%%%%%%%%%%%%%%%
\subsection{Greek Letters Used for Angles}
\label{subsec:greek-angle-letters}
%%%%%%%%%%%%%%%%%%%%%%%%%%%%%%%%%%%%%%%%%%%%%%%%%%%%%%%%%%%%

Angles are commonly represented by

\[
\boxed{
\alpha,
\quad
\beta,
\quad
\gamma,
\quad
\theta,
\quad
\phi.
}
\]

For example,

\[
\boxed{
\sin\theta,
\qquad
\cos\theta.
}
\]

No Greek letter is inherently an angle; this is simply a common convention.

%%%%%%%%%%%%%%%%%%%%%%%%%%%%%%%%%%%%%%%%%%%%%%%%%%%%%%%%%%%%
\subsection{Greek Letters in Calculus and Analysis}
\label{subsec:greek-analysis-uses}
%%%%%%%%%%%%%%%%%%%%%%%%%%%%%%%%%%%%%%%%%%%%%%%%%%%%%%%%%%%%

Common symbols include:

\[
\boxed{
\epsilon,\delta
}
\]

for limit definitions,

\[
\boxed{
\lambda
}
\]

for eigenvalues and parameters,

\[
\boxed{
\mu
}
\]

for measures,

and

\[
\boxed{
\omega
}
\]

for differential forms or moduli of continuity.

For example,

\[
\boxed{
|x-a|<\delta
\Rightarrow
|f(x)-L|<\varepsilon.
}
\]

%%%%%%%%%%%%%%%%%%%%%%%%%%%%%%%%%%%%%%%%%%%%%%%%%%%%%%%%%%%%
\subsection{Greek Letters in Linear Algebra}
\label{subsec:greek-linear-algebra-uses}
%%%%%%%%%%%%%%%%%%%%%%%%%%%%%%%%%%%%%%%%%%%%%%%%%%%%%%%%%%%%

Eigenvalues are frequently denoted

\[
\boxed{
\lambda_1,
\lambda_2,
\ldots,
\lambda_n.
}
\]

Thus

\[
\boxed{
A\mathbf{v}_i
=
\lambda_i\mathbf{v}_i.
}
\]

Singular values are often written

\[
\boxed{
\sigma_1,
\sigma_2,
\ldots.
}
\]

%%%%%%%%%%%%%%%%%%%%%%%%%%%%%%%%%%%%%%%%%%%%%%%%%%%%%%%%%%%%
\subsection{Greek Letters in Probability}
\label{subsec:greek-probability-uses}
%%%%%%%%%%%%%%%%%%%%%%%%%%%%%%%%%%%%%%%%%%%%%%%%%%%%%%%%%%%%

Common notation includes:

\[
\boxed{
\Omega
}
\]

for the sample space,

\[
\boxed{
\mu
}
\]

for a mean or measure,

\[
\boxed{
\sigma^2
}
\]

for variance,

\[
\boxed{
\lambda
}
\]

for Poisson or exponential rates,

and

\[
\boxed{
\rho
}
\]

for correlation.

%%%%%%%%%%%%%%%%%%%%%%%%%%%%%%%%%%%%%%%%%%%%%%%%%%%%%%%%%%%%
\subsection{Greek Letters in Statistics}
\label{subsec:greek-statistics-uses}
%%%%%%%%%%%%%%%%%%%%%%%%%%%%%%%%%%%%%%%%%%%%%%%%%%%%%%%%%%%%

Common statistical notation includes

\[
\boxed{
\mu
=
\text{population mean},
}
\]

\[
\boxed{
\sigma
=
\text{population standard deviation},
}
\]

\[
\boxed{
\alpha
=
\text{significance level},
}
\]

\[
\boxed{
\beta
=
\text{regression coefficient or Type II error probability},
}
\]

and

\[
\boxed{
\theta
=
\text{generic parameter}.
}
\]

%%%%%%%%%%%%%%%%%%%%%%%%%%%%%%%%%%%%%%%%%%%%%%%%%%%%%%%%%%%%
\subsection{Greek Letters in Differential Equations}
\label{subsec:greek-differential-equation-uses}
%%%%%%%%%%%%%%%%%%%%%%%%%%%%%%%%%%%%%%%%%%%%%%%%%%%%%%%%%%%%

Parameters may be denoted by

\[
\boxed{
\lambda,
\quad
\mu,
\quad
\gamma,
\quad
\omega.
}
\]

For example,

\[
\boxed{
y''+\omega^2y=0
}
\]

describes a harmonic oscillator.

The parameter \(\omega\) is angular frequency.

%%%%%%%%%%%%%%%%%%%%%%%%%%%%%%%%%%%%%%%%%%%%%%%%%%%%%%%%%%%%
\subsection{Greek Letters in Optimization}
\label{subsec:greek-optimization-uses}
%%%%%%%%%%%%%%%%%%%%%%%%%%%%%%%%%%%%%%%%%%%%%%%%%%%%%%%%%%%%

Lagrange multipliers are often written

\[
\boxed{
\lambda_i
}
\]

and

\[
\boxed{
\mu_j.
}
\]

A constrained Lagrangian may be

\[
\boxed{
\mathcal{L}
=
f(x)
+
\sum_i
\lambda_i h_i(x)
+
\sum_j
\mu_j g_j(x).
}
\]

%%%%%%%%%%%%%%%%%%%%%%%%%%%%%%%%%%%%%%%%%%%%%%%%%%%%%%%%%%%%
\subsection{Greek Letters in Geometry}
\label{subsec:greek-geometry-uses}
%%%%%%%%%%%%%%%%%%%%%%%%%%%%%%%%%%%%%%%%%%%%%%%%%%%%%%%%%%%%

Angles are often

\[
\boxed{
\alpha,\beta,\gamma,\theta,\phi.
}
\]

Curvature may be written

\[
\boxed{
\kappa.
}
\]

A metric or geometric parameter may be written using many other Greek
letters depending on the subject.

%%%%%%%%%%%%%%%%%%%%%%%%%%%%%%%%%%%%%%%%%%%%%%%%%%%%%%%%%%%%
\subsection{Greek Letters in Topology}
\label{subsec:greek-topology-uses}
%%%%%%%%%%%%%%%%%%%%%%%%%%%%%%%%%%%%%%%%%%%%%%%%%%%%%%%%%%%%

Topology commonly uses Greek letters for maps and indices.

For example,

\[
\boxed{
\phi:X\to Y
}
\]

may denote a homeomorphism.

A topology itself is sometimes written

\[
\boxed{
\tau.
}
\]

Thus a topological space may be denoted

\[
\boxed{
(X,\tau).
}
\]

%%%%%%%%%%%%%%%%%%%%%%%%%%%%%%%%%%%%%%%%%%%%%%%%%%%%%%%%%%%%
\subsection{Greek Letters in Abstract Algebra}
\label{subsec:greek-algebra-uses}
%%%%%%%%%%%%%%%%%%%%%%%%%%%%%%%%%%%%%%%%%%%%%%%%%%%%%%%%%%%%

Homomorphisms are frequently written

\[
\boxed{
\phi,
\quad
\psi.
}
\]

For example,

\[
\boxed{
\phi:G\to H.
}
\]

Representations, characters, roots, and weights also commonly use Greek
letters.

%%%%%%%%%%%%%%%%%%%%%%%%%%%%%%%%%%%%%%%%%%%%%%%%%%%%%%%%%%%%
\subsection{Greek Letters in Number Theory}
\label{subsec:greek-number-theory-uses}
%%%%%%%%%%%%%%%%%%%%%%%%%%%%%%%%%%%%%%%%%%%%%%%%%%%%%%%%%%%%

Common examples include

\[
\boxed{
\phi(n)
}
\]

for Euler's totient function,

\[
\boxed{
\mu(n)
}
\]

for the M\"obius function,

and

\[
\boxed{
\zeta(s)
}
\]

for the Riemann zeta function.

The same symbol may represent something completely different in another
subject.

%%%%%%%%%%%%%%%%%%%%%%%%%%%%%%%%%%%%%%%%%%%%%%%%%%%%%%%%%%%%
\subsection{Greek Letters in Physics}
\label{subsec:greek-physics-uses}
%%%%%%%%%%%%%%%%%%%%%%%%%%%%%%%%%%%%%%%%%%%%%%%%%%%%%%%%%%%%

Physics uses Greek notation extensively.

Common examples include

\[
\boxed{
\rho
=
\text{density},
}
\]

\[
\boxed{
\lambda
=
\text{wavelength},
}
\]

\[
\boxed{
\omega
=
\text{angular frequency},
}
\]

\[
\boxed{
\mu
=
\text{coefficient or permeability},
}
\]

and

\[
\boxed{
\psi
=
\text{wave function}.
}
\]

%%%%%%%%%%%%%%%%%%%%%%%%%%%%%%%%%%%%%%%%%%%%%%%%%%%%%%%%%%%%
\subsection{Greek Letters in Continuum Mechanics}
\label{subsec:greek-continuum-uses}
%%%%%%%%%%%%%%%%%%%%%%%%%%%%%%%%%%%%%%%%%%%%%%%%%%%%%%%%%%%%

Typical notation includes

\[
\boxed{
\rho
=
\text{density},
}
\]

\[
\boxed{
\boldsymbol{\sigma}
=
\text{stress tensor},
}
\]

\[
\boxed{
\boldsymbol{\varepsilon}
=
\text{strain tensor},
}
\]

and

\[
\boxed{
\mu,\lambda
}
\]

for Lam\'e-type material parameters.

%%%%%%%%%%%%%%%%%%%%%%%%%%%%%%%%%%%%%%%%%%%%%%%%%%%%%%%%%%%%
\subsection{Greek Letters in Fluid Dynamics}
\label{subsec:greek-fluid-uses}
%%%%%%%%%%%%%%%%%%%%%%%%%%%%%%%%%%%%%%%%%%%%%%%%%%%%%%%%%%%%

Common examples are

\[
\boxed{
\rho
=
\text{density},
}
\]

\[
\boxed{
\mu
=
\text{dynamic viscosity},
}
\]

\[
\boxed{
\nu
=
\text{kinematic viscosity},
}
\]

\[
\boxed{
\boldsymbol{\omega}
=
\text{vorticity},
}
\]

and

\[
\boxed{
\Gamma
=
\text{circulation}.
}
\]

%%%%%%%%%%%%%%%%%%%%%%%%%%%%%%%%%%%%%%%%%%%%%%%%%%%%%%%%%%%%
\subsection{Greek Letters in Control Theory}
\label{subsec:greek-control-uses}
%%%%%%%%%%%%%%%%%%%%%%%%%%%%%%%%%%%%%%%%%%%%%%%%%%%%%%%%%%%%

Control theory often uses

\[
\boxed{
\lambda_i
}
\]

for system eigenvalues,

\[
\boxed{
\omega_n
}
\]

for natural frequency,

and

\[
\boxed{
\zeta
}
\]

for damping ratio.

For a second-order system,

\[
\boxed{
s^2
+
2\zeta\omega_ns
+
\omega_n^2.
}
\]

%%%%%%%%%%%%%%%%%%%%%%%%%%%%%%%%%%%%%%%%%%%%%%%%%%%%%%%%%%%%
\subsection{Greek Letters in Finance}
\label{subsec:greek-finance-uses}
%%%%%%%%%%%%%%%%%%%%%%%%%%%%%%%%%%%%%%%%%%%%%%%%%%%%%%%%%%%%

Mathematical finance uses Greek letters heavily for option sensitivities.

The option Greeks include

\[
\boxed{
\Delta
}
\]

for delta,

\[
\boxed{
\Gamma
}
\]

for gamma,

\[
\boxed{
\Theta
}
\]

for theta,

and

\[
\boxed{
\rho
}
\]

for rho.

Volatility is typically

\[
\boxed{
\sigma.
}
\]

%%%%%%%%%%%%%%%%%%%%%%%%%%%%%%%%%%%%%%%%%%%%%%%%%%%%%%%%%%%%
\subsection{Greek Letters in Economics}
\label{subsec:greek-economics-uses}
%%%%%%%%%%%%%%%%%%%%%%%%%%%%%%%%%%%%%%%%%%%%%%%%%%%%%%%%%%%%

Economic notation often includes

\[
\boxed{
\alpha,\beta
}
\]

for elasticities, production exponents, preferences, or discount parameters,

\[
\boxed{
\delta
}
\]

for depreciation,

and

\[
\boxed{
\lambda
}
\]

for multipliers or shadow prices.

For Cobb--Douglas production,

\[
\boxed{
Y
=
AK^\alpha L^\beta.
}
\]

%%%%%%%%%%%%%%%%%%%%%%%%%%%%%%%%%%%%%%%%%%%%%%%%%%%%%%%%%%%%
\subsection{Greek Letters as Distribution Parameters}
\label{subsec:greek-distribution-parameters}
%%%%%%%%%%%%%%%%%%%%%%%%%%%%%%%%%%%%%%%%%%%%%%%%%%%%%%%%%%%%

Probability distributions frequently use Greek parameters.

Examples include:

\[
\boxed{
\lambda
}
\]

for Poisson and exponential distributions,

\[
\boxed{
\mu,\sigma
}
\]

for normal distributions,

and

\[
\boxed{
\alpha,\beta
}
\]

for beta and gamma-family distributions.

Thus

\[
\boxed{
X
\sim
N(\mu,\sigma^2).
}
\]

%%%%%%%%%%%%%%%%%%%%%%%%%%%%%%%%%%%%%%%%%%%%%%%%%%%%%%%%%%%%
\subsection{Greek Letters as Indices}
\label{subsec:greek-indices}
%%%%%%%%%%%%%%%%%%%%%%%%%%%%%%%%%%%%%%%%%%%%%%%%%%%%%%%%%%%%

In advanced mathematics and physics, Greek letters may serve as tensor
indices.

For example,

\[
\boxed{
T_{\mu\nu}.
}
\]

Latin and Greek indices may intentionally refer to different coordinate
ranges.

Whenever this convention is used, the index ranges should be stated.

%%%%%%%%%%%%%%%%%%%%%%%%%%%%%%%%%%%%%%%%%%%%%%%%%%%%%%%%%%%%
\subsection{Greek Letters as Maps}
\label{subsec:greek-maps}
%%%%%%%%%%%%%%%%%%%%%%%%%%%%%%%%%%%%%%%%%%%%%%%%%%%%%%%%%%%%

Maps are commonly named

\[
\boxed{
\phi,
\psi,
\varphi.
}
\]

For example,

\[
\boxed{
\phi:G\to H
}
\]

may denote a group homomorphism.

Similarly,

\[
\boxed{
\psi:X\to\mathbb{C}
}
\]

may represent a function or state.

%%%%%%%%%%%%%%%%%%%%%%%%%%%%%%%%%%%%%%%%%%%%%%%%%%%%%%%%%%%%
\subsection{Greek Letters as Measures}
\label{subsec:greek-measures}
%%%%%%%%%%%%%%%%%%%%%%%%%%%%%%%%%%%%%%%%%%%%%%%%%%%%%%%%%%%%

Measures are often written

\[
\boxed{
\mu,
\quad
\nu.
}
\]

For example,

\[
\boxed{
\mu(A)
}
\]

may denote the measure of set \(A\).

A statement such as

\[
\boxed{
\mu\ll\nu
}
\]

means that \(\mu\) is absolutely continuous with respect to \(\nu\).

%%%%%%%%%%%%%%%%%%%%%%%%%%%%%%%%%%%%%%%%%%%%%%%%%%%%%%%%%%%%
\subsection{Greek Letters as Eigenvalues}
\label{subsec:greek-eigenvalues}
%%%%%%%%%%%%%%%%%%%%%%%%%%%%%%%%%%%%%%%%%%%%%%%%%%%%%%%%%%%%

The most common eigenvalue symbol is

\[
\boxed{
\lambda.
}
\]

For several eigenvalues,

\[
\boxed{
\lambda_1,
\lambda_2,
\ldots,
\lambda_n.
}
\]

In spectral problems one may also encounter

\[
\boxed{
\mu,
\quad
\sigma,
\quad
\omega.
}
\]

%%%%%%%%%%%%%%%%%%%%%%%%%%%%%%%%%%%%%%%%%%%%%%%%%%%%%%%%%%%%
\subsection{Greek Letters as Parameters}
\label{subsec:greek-parameters}
%%%%%%%%%%%%%%%%%%%%%%%%%%%%%%%%%%%%%%%%%%%%%%%%%%%%%%%%%%%%

A generic model may be written

\[
\boxed{
f(x;\theta).
}
\]

Here

\[
\boxed{
\theta
}
\]

represents an unspecified parameter or parameter vector.

Other common generic parameter letters include

\[
\boxed{
\alpha,\beta,\lambda,\mu.
}
\]

%%%%%%%%%%%%%%%%%%%%%%%%%%%%%%%%%%%%%%%%%%%%%%%%%%%%%%%%%%%%
\subsection{Greek Letters as Random Variables}
\label{subsec:greek-random-variables}
%%%%%%%%%%%%%%%%%%%%%%%%%%%%%%%%%%%%%%%%%%%%%%%%%%%%%%%%%%%%

Random variables may also be represented by Greek letters:

\[
\boxed{
\xi,
\quad
\eta,
\quad
\zeta.
}
\]

For example,

\[
\boxed{
\xi\sim N(0,1).
}
\]

Capital Latin letters remain more common in elementary probability.

%%%%%%%%%%%%%%%%%%%%%%%%%%%%%%%%%%%%%%%%%%%%%%%%%%%%%%%%%%%%
\subsection{Greek Letters and Dimensions}
\label{subsec:greek-dimensions}
%%%%%%%%%%%%%%%%%%%%%%%%%%%%%%%%%%%%%%%%%%%%%%%%%%%%%%%%%%%%

A Greek symbol does not determine physical units.

For example,

\[
\lambda
\]

could represent:

\[
\boxed{
\text{a dimensionless eigenvalue}
}
\]

or

\[
\boxed{
\text{a wavelength measured in meters}.
}
\]

Units must therefore be inferred from the mathematical definition, not the
symbol itself.

%%%%%%%%%%%%%%%%%%%%%%%%%%%%%%%%%%%%%%%%%%%%%%%%%%%%%%%%%%%%
\subsection{Commonly Confused Greek Symbols}
\label{subsec:confused-greek-symbols}
%%%%%%%%%%%%%%%%%%%%%%%%%%%%%%%%%%%%%%%%%%%%%%%%%%%%%%%%%%%%

Several Greek letters can look similar to Latin letters or other symbols.

Examples include

\[
\boxed{
\nu
\quad\text{and}\quad
v,
}
\]

\[
\boxed{
\rho
\quad\text{and}\quad
p,
}
\]

\[
\boxed{
\chi
\quad\text{and}\quad
x,
}
\]

\[
\boxed{
\eta
\quad\text{and}\quad
n,
}
\]

and

\[
\boxed{
\omega
\quad\text{and}\quad
w.
}
\]

Careful handwriting and typesetting reduce ambiguity.

%%%%%%%%%%%%%%%%%%%%%%%%%%%%%%%%%%%%%%%%%%%%%%%%%%%%%%%%%%%%
\subsection{Epsilon Versus Membership}
\label{subsec:greek-epsilon-versus-membership}
%%%%%%%%%%%%%%%%%%%%%%%%%%%%%%%%%%%%%%%%%%%%%%%%%%%%%%%%%%%%

The Greek letters

\[
\boxed{
\epsilon,
\quad
\varepsilon
}
\]

must not be confused with the membership symbol

\[
\boxed{
\in.
}
\]

For example,

\[
\boxed{
\varepsilon\in\mathbb{R}
}
\]

contains both symbols with different meanings.

%%%%%%%%%%%%%%%%%%%%%%%%%%%%%%%%%%%%%%%%%%%%%%%%%%%%%%%%%%%%
\subsection{Nu Versus Latin V}
\label{subsec:nu-versus-v}
%%%%%%%%%%%%%%%%%%%%%%%%%%%%%%%%%%%%%%%%%%%%%%%%%%%%%%%%%%%%

The Greek letter

\[
\boxed{
\nu
}
\]

can look similar to the Latin letter

\[
\boxed{
v.
}
\]

This is especially relevant in fluid mechanics, where

\[
\nu
\]

may represent kinematic viscosity while

\[
v
\]

may represent a velocity component.

%%%%%%%%%%%%%%%%%%%%%%%%%%%%%%%%%%%%%%%%%%%%%%%%%%%%%%%%%%%%
\subsection{Rho Versus Latin P}
\label{subsec:rho-versus-p}
%%%%%%%%%%%%%%%%%%%%%%%%%%%%%%%%%%%%%%%%%%%%%%%%%%%%%%%%%%%%

The Greek letter

\[
\boxed{
\rho
}
\]

can resemble the Latin letter \(p\) in handwriting.

In fluid mechanics,

\[
\boxed{
\rho
=
\text{density}
}
\]

while

\[
\boxed{
p
=
\text{pressure}.
}
\]

These quantities are physically and mathematically distinct.

%%%%%%%%%%%%%%%%%%%%%%%%%%%%%%%%%%%%%%%%%%%%%%%%%%%%%%%%%%%%
\subsection{Chi Versus Latin X}
\label{subsec:chi-versus-x}
%%%%%%%%%%%%%%%%%%%%%%%%%%%%%%%%%%%%%%%%%%%%%%%%%%%%%%%%%%%%

The Greek letter

\[
\boxed{
\chi
}
\]

resembles \(x\).

For example,

\[
\boxed{
\chi^2
}
\]

is read ``chi squared,'' not ``x squared.''

%%%%%%%%%%%%%%%%%%%%%%%%%%%%%%%%%%%%%%%%%%%%%%%%%%%%%%%%%%%%
\subsection{Phi Variants}
\label{subsec:phi-variants}
%%%%%%%%%%%%%%%%%%%%%%%%%%%%%%%%%%%%%%%%%%%%%%%%%%%%%%%%%%%%

The symbols

\[
\boxed{
\phi
}
\]

and

\[
\boxed{
\varphi
}
\]

are alternate forms of the same Greek letter.

A book may use them simultaneously for different objects.

When that occurs, their meanings must be explicitly defined.

%%%%%%%%%%%%%%%%%%%%%%%%%%%%%%%%%%%%%%%%%%%%%%%%%%%%%%%%%%%%
\subsection{Epsilon Variants}
\label{subsec:epsilon-variants}
%%%%%%%%%%%%%%%%%%%%%%%%%%%%%%%%%%%%%%%%%%%%%%%%%%%%%%%%%%%%

Similarly,

\[
\boxed{
\epsilon
}
\]

and

\[
\boxed{
\varepsilon
}
\]

are alternate forms.

Some texts prefer \(\varepsilon\) for formal analysis because it is visually
distinct from the membership symbol \(\in\).

%%%%%%%%%%%%%%%%%%%%%%%%%%%%%%%%%%%%%%%%%%%%%%%%%%%%%%%%%%%%
\subsection{Theta Variants}
\label{subsec:theta-variants}
%%%%%%%%%%%%%%%%%%%%%%%%%%%%%%%%%%%%%%%%%%%%%%%%%%%%%%%%%%%%

The symbols

\[
\boxed{
\theta
}
\]

and

\[
\boxed{
\vartheta
}
\]

are alternate theta forms.

The standard form

\[
\theta
\]

is especially common for angles.

%%%%%%%%%%%%%%%%%%%%%%%%%%%%%%%%%%%%%%%%%%%%%%%%%%%%%%%%%%%%
\subsection{Sigma Variants}
\label{subsec:sigma-variants}
%%%%%%%%%%%%%%%%%%%%%%%%%%%%%%%%%%%%%%%%%%%%%%%%%%%%%%%%%%%%

The symbols

\[
\boxed{
\sigma
}
\]

and

\[
\boxed{
\varsigma
}
\]

are lowercase sigma variants.

The form

\[
\sigma
\]

is overwhelmingly more common in general mathematical notation.

%%%%%%%%%%%%%%%%%%%%%%%%%%%%%%%%%%%%%%%%%%%%%%%%%%%%%%%%%%%%
\subsection{Pi Variants}
\label{subsec:pi-variants}
%%%%%%%%%%%%%%%%%%%%%%%%%%%%%%%%%%%%%%%%%%%%%%%%%%%%%%%%%%%%

The usual circle constant is

\[
\boxed{
\pi.
}
\]

The alternate symbol

\[
\boxed{
\varpi
}
\]

appears occasionally in advanced mathematical contexts.

It should not normally replace \(\pi\) for the circle constant unless a
specific notation convention requires it.

%%%%%%%%%%%%%%%%%%%%%%%%%%%%%%%%%%%%%%%%%%%%%%%%%%%%%%%%%%%%
\subsection{Uppercase Greek Operators}
\label{subsec:uppercase-greek-operators}
%%%%%%%%%%%%%%%%%%%%%%%%%%%%%%%%%%%%%%%%%%%%%%%%%%%%%%%%%%%%

Several uppercase Greek letters have major operator roles.

The summation operator is

\[
\boxed{
\Sigma
\longrightarrow
\sum.
}
\]

The product operator is

\[
\boxed{
\Pi
\longrightarrow
\prod.
}
\]

The finite-change or Laplacian symbol is

\[
\boxed{
\Delta.
}
\]

These operator forms may differ stylistically from ordinary uppercase Greek
letters in typeset mathematics.

%%%%%%%%%%%%%%%%%%%%%%%%%%%%%%%%%%%%%%%%%%%%%%%%%%%%%%%%%%%%
\subsection{Greek Letters in LaTeX}
\label{subsec:greek-letters-latex}
%%%%%%%%%%%%%%%%%%%%%%%%%%%%%%%%%%%%%%%%%%%%%%%%%%%%%%%%%%%%

Greek letters are typed inside math mode.

For example,

\begin{verbatim}
\[
\alpha + \beta = \gamma
\]
\end{verbatim}

produces

\[
\boxed{
\alpha+\beta=\gamma.
}
\]

Uppercase Greek letters with distinct forms use capitalized commands:

\begin{verbatim}
\[
\Gamma,\Delta,\Theta,\Lambda,\Xi,\Pi,\Sigma,\Upsilon,\Phi,\Psi,\Omega
\]
\end{verbatim}

%%%%%%%%%%%%%%%%%%%%%%%%%%%%%%%%%%%%%%%%%%%%%%%%%%%%%%%%%%%%
\subsection{Lowercase LaTeX Reference}
\label{subsec:lowercase-greek-latex-reference}
%%%%%%%%%%%%%%%%%%%%%%%%%%%%%%%%%%%%%%%%%%%%%%%%%%%%%%%%%%%%

\begin{center}
\begin{tabular}{c|l}
\textbf{Symbol} & \textbf{LaTeX} \\ \hline
\(\alpha\) & \verb|\alpha| \\
\(\beta\) & \verb|\beta| \\
\(\gamma\) & \verb|\gamma| \\
\(\delta\) & \verb|\delta| \\
\(\epsilon\) & \verb|\epsilon| \\
\(\varepsilon\) & \verb|\varepsilon| \\
\(\zeta\) & \verb|\zeta| \\
\(\eta\) & \verb|\eta| \\
\(\theta\) & \verb|\theta| \\
\(\vartheta\) & \verb|\vartheta| \\
\(\iota\) & \verb|\iota| \\
\(\kappa\) & \verb|\kappa| \\
\(\lambda\) & \verb|\lambda| \\
\(\mu\) & \verb|\mu| \\
\(\nu\) & \verb|\nu| \\
\(\xi\) & \verb|\xi| \\
\(\pi\) & \verb|\pi| \\
\(\varpi\) & \verb|\varpi| \\
\(\rho\) & \verb|\rho| \\
\(\varrho\) & \verb|\varrho| \\
\(\sigma\) & \verb|\sigma| \\
\(\varsigma\) & \verb|\varsigma| \\
\(\tau\) & \verb|\tau| \\
\(\upsilon\) & \verb|\upsilon| \\
\(\phi\) & \verb|\phi| \\
\(\varphi\) & \verb|\varphi| \\
\(\chi\) & \verb|\chi| \\
\(\psi\) & \verb|\psi| \\
\(\omega\) & \verb|\omega|
\end{tabular}
\end{center}

%%%%%%%%%%%%%%%%%%%%%%%%%%%%%%%%%%%%%%%%%%%%%%%%%%%%%%%%%%%%
\subsection{Uppercase LaTeX Reference}
\label{subsec:uppercase-greek-latex-reference}
%%%%%%%%%%%%%%%%%%%%%%%%%%%%%%%%%%%%%%%%%%%%%%%%%%%%%%%%%%%%

\begin{center}
\begin{tabular}{c|l}
\textbf{Symbol} & \textbf{LaTeX} \\ \hline
\(\Gamma\) & \verb|\Gamma| \\
\(\Delta\) & \verb|\Delta| \\
\(\Theta\) & \verb|\Theta| \\
\(\Lambda\) & \verb|\Lambda| \\
\(\Xi\) & \verb|\Xi| \\
\(\Pi\) & \verb|\Pi| \\
\(\Sigma\) & \verb|\Sigma| \\
\(\Upsilon\) & \verb|\Upsilon| \\
\(\Phi\) & \verb|\Phi| \\
\(\Psi\) & \verb|\Psi| \\
\(\Omega\) & \verb|\Omega|
\end{tabular}
\end{center}

%%%%%%%%%%%%%%%%%%%%%%%%%%%%%%%%%%%%%%%%%%%%%%%%%%%%%%%%%%%%
\subsection{Bold Greek Letters}
\label{subsec:bold-greek-letters}
%%%%%%%%%%%%%%%%%%%%%%%%%%%%%%%%%%%%%%%%%%%%%%%%%%%%%%%%%%%%

Greek symbols may represent vectors or tensors.

With appropriate LaTeX support, bold Greek may be written using

\begin{verbatim}
\boldsymbol{\sigma}
\end{verbatim}

to produce

\[
\boxed{
\boldsymbol{\sigma}.
}
\]

Similarly,

\[
\boxed{
\boldsymbol{\theta}
}
\]

may represent a parameter vector.

This is preferable to changing the symbol itself when boldface already
indicates vector or tensor quantities elsewhere in the text.

%%%%%%%%%%%%%%%%%%%%%%%%%%%%%%%%%%%%%%%%%%%%%%%%%%%%%%%%%%%%
\subsection{Subscripts and Superscripts}
\label{subsec:greek-subscripts-superscripts}
%%%%%%%%%%%%%%%%%%%%%%%%%%%%%%%%%%%%%%%%%%%%%%%%%%%%%%%%%%%%

Greek letters combine with indices normally:

\[
\boxed{
\lambda_i,
}
\]

\[
\boxed{
\sigma_{ij},
}
\]

\[
\boxed{
\theta^{(k)},
}
\]

and

\[
\boxed{
\omega_n.
}
\]

For example,

\[
\boxed{
\lambda_1\leq\lambda_2\leq\cdots\leq\lambda_n.
}
\]

%%%%%%%%%%%%%%%%%%%%%%%%%%%%%%%%%%%%%%%%%%%%%%%%%%%%%%%%%%%%
\subsection{Greek Letters With Hats, Bars, and Tildes}
\label{subsec:decorated-greek-letters}
%%%%%%%%%%%%%%%%%%%%%%%%%%%%%%%%%%%%%%%%%%%%%%%%%%%%%%%%%%%%

Greek symbols may be decorated like Latin symbols.

Examples include

\[
\boxed{
\hat{\theta}
}
\]

for an estimate,

\[
\boxed{
\bar{\theta}
}
\]

for an average or reference value,

and

\[
\boxed{
\tilde{\theta}
}
\]

for a transformed or approximate parameter.

The decoration itself has no universal meaning.

%%%%%%%%%%%%%%%%%%%%%%%%%%%%%%%%%%%%%%%%%%%%%%%%%%%%%%%%%%%%
\subsection{Greek Letters in Subscripts}
\label{subsec:greek-in-subscripts}
%%%%%%%%%%%%%%%%%%%%%%%%%%%%%%%%%%%%%%%%%%%%%%%%%%%%%%%%%%%%

Greek indices are common in tensor notation:

\[
\boxed{
T_{\mu\nu}.
}
\]

If

\[
\mu,\nu
=
0,1,2,3,
\]

then the notation may refer to spacetime coordinates.

In another text, the same symbols may have completely different ranges.

Always define index conventions.

%%%%%%%%%%%%%%%%%%%%%%%%%%%%%%%%%%%%%%%%%%%%%%%%%%%%%%%%%%%%
\subsection{Symbol Choice Guidelines}
\label{subsec:greek-symbol-choice-guidelines}
%%%%%%%%%%%%%%%%%%%%%%%%%%%%%%%%%%%%%%%%%%%%%%%%%%%%%%%%%%%%

When choosing Greek symbols:

\begin{enumerate}
    \item follow established conventions when they improve readability;
    \item avoid giving the same symbol multiple meanings in one local context;
    \item avoid visually similar symbols when possible;
    \item define alternate forms explicitly;
    \item preserve distinctions between scalars, vectors, and tensors;
    \item state units when the quantity is physical;
    \item state index ranges;
    \item distinguish parameters from variables when helpful;
    \item do not assume a conventional meaning is universal.
\end{enumerate}

%%%%%%%%%%%%%%%%%%%%%%%%%%%%%%%%%%%%%%%%%%%%%%%%%%%%%%%%%%%%
\subsection{Common Errors}
\label{subsec:greek-alphabet-common-errors}
%%%%%%%%%%%%%%%%%%%%%%%%%%%%%%%%%%%%%%%%%%%%%%%%%%%%%%%%%%%%

\subsubsection{Assuming Greek Letters Have Fixed Meanings}

The symbol

\[
\lambda
\]

does not inherently mean eigenvalue.

It only has that meaning after the notation is defined or established by
context.

%%%%%%%%%%%%%%%%%%%%%%%%%%%%%%%%%%%%%%%%%%%%%%%%%%%%%%%%%%%%
\subsubsection{Confusing Epsilon and Membership}

Do not confuse

\[
\varepsilon
\]

with

\[
\in.
\]

For example,

\[
\boxed{
\varepsilon\in(0,1)
}
\]

uses epsilon as a variable and \(\in\) as set membership.

%%%%%%%%%%%%%%%%%%%%%%%%%%%%%%%%%%%%%%%%%%%%%%%%%%%%%%%%%%%%
\subsubsection{Confusing Nu and Velocity}

In printed mathematics,

\[
\nu
\]

and

\[
v
\]

can be similar.

Use careful typesetting when both appear in the same formula.

%%%%%%%%%%%%%%%%%%%%%%%%%%%%%%%%%%%%%%%%%%%%%%%%%%%%%%%%%%%%
\subsubsection{Confusing Rho and Pressure}

In fluid mechanics,

\[
\rho
\]

often denotes density while

\[
p
\]

denotes pressure.

Poor handwriting can make them difficult to distinguish.

%%%%%%%%%%%%%%%%%%%%%%%%%%%%%%%%%%%%%%%%%%%%%%%%%%%%%%%%%%%%
\subsubsection{Confusing Chi and \(x\)}

The symbol

\[
\chi
\]

is a Greek letter.

The symbol

\[
x
\]

is a Latin variable.

For example,

\[
\chi^2
\]

is pronounced ``chi squared.''

%%%%%%%%%%%%%%%%%%%%%%%%%%%%%%%%%%%%%%%%%%%%%%%%%%%%%%%%%%%%
\subsubsection{Assuming \(\phi\) and \(\varphi\) Are Different Letters}

They are typographic forms of the same Greek letter.

A text may assign different meanings to them, but the distinction is local
notation rather than a difference in the alphabet.

%%%%%%%%%%%%%%%%%%%%%%%%%%%%%%%%%%%%%%%%%%%%%%%%%%%%%%%%%%%%
\subsubsection{Using Uppercase Commands That Do Not Exist}

Not every uppercase Greek letter has a distinct LaTeX command.

For example, uppercase alpha looks like Latin \(A\), so one normally writes

\[
\boxed{
A
}
\]

rather than attempting a separate Greek-alpha command.

%%%%%%%%%%%%%%%%%%%%%%%%%%%%%%%%%%%%%%%%%%%%%%%%%%%%%%%%%%%%
\subsubsection{Using Greek Letters Outside Math Mode}

In LaTeX, commands such as

\begin{verbatim}
\alpha
\end{verbatim}

must normally be used inside mathematical mode.

For example,

\begin{verbatim}
$\alpha$
\end{verbatim}

or

\begin{verbatim}
\[
\alpha
\]
\end{verbatim}

is appropriate.

%%%%%%%%%%%%%%%%%%%%%%%%%%%%%%%%%%%%%%%%%%%%%%%%%%%%%%%%%%%%
\subsubsection{Changing Notation Midway Through a Derivation}

If

\[
\sigma
\]

represents standard deviation at the beginning of a calculation, do not later
reuse it for stress or another unrelated quantity without explicitly
redefining it.

%%%%%%%%%%%%%%%%%%%%%%%%%%%%%%%%%%%%%%%%%%%%%%%%%%%%%%%%%%%%
\subsubsection{Ignoring Boldface Distinctions}

If

\[
\sigma
\]

is a scalar while

\[
\boldsymbol{\sigma}
\]

is a tensor, removing the boldface changes the mathematical object.

%%%%%%%%%%%%%%%%%%%%%%%%%%%%%%%%%%%%%%%%%%%%%%%%%%%%%%%%%%%%
\subsection{Quick Greek Alphabet Reference}
\label{subsec:quick-greek-reference}
%%%%%%%%%%%%%%%%%%%%%%%%%%%%%%%%%%%%%%%%%%%%%%%%%%%%%%%%%%%%

\begin{center}
\begin{tabular}{c|c|c}
\textbf{Name}
&
\textbf{Lowercase}
&
\textbf{Uppercase}
\\ \hline

Alpha & \(\alpha\) & \(A\) \\
Beta & \(\beta\) & \(B\) \\
Gamma & \(\gamma\) & \(\Gamma\) \\
Delta & \(\delta\) & \(\Delta\) \\
Epsilon & \(\epsilon,\varepsilon\) & \(E\) \\
Zeta & \(\zeta\) & \(Z\) \\
Eta & \(\eta\) & \(H\) \\
Theta & \(\theta,\vartheta\) & \(\Theta\) \\
Iota & \(\iota\) & \(I\) \\
Kappa & \(\kappa\) & \(K\) \\
Lambda & \(\lambda\) & \(\Lambda\) \\
Mu & \(\mu\) & \(M\) \\
Nu & \(\nu\) & \(N\) \\
Xi & \(\xi\) & \(\Xi\) \\
Omicron & \(o\) & \(O\) \\
Pi & \(\pi,\varpi\) & \(\Pi\) \\
Rho & \(\rho,\varrho\) & \(P\) \\
Sigma & \(\sigma,\varsigma\) & \(\Sigma\) \\
Tau & \(\tau\) & \(T\) \\
Upsilon & \(\upsilon\) & \(\Upsilon\) \\
Phi & \(\phi,\varphi\) & \(\Phi\) \\
Chi & \(\chi\) & \(X\) \\
Psi & \(\psi\) & \(\Psi\) \\
Omega & \(\omega\) & \(\Omega\)
\end{tabular}
\end{center}

%%%%%%%%%%%%%%%%%%%%%%%%%%%%%%%%%%%%%%%%%%%%%%%%%%%%%%%%%%%%
\subsection{Quick Common-Use Reference}
\label{subsec:quick-greek-use-reference}
%%%%%%%%%%%%%%%%%%%%%%%%%%%%%%%%%%%%%%%%%%%%%%%%%%%%%%%%%%%%

\begin{center}
\begin{tabular}{c|l}
\textbf{Symbol} & \textbf{Common Uses} \\ \hline
\(\alpha\) & angle, significance level, parameter \\
\(\beta\) & regression coefficient, rate, parameter \\
\(\gamma\) & rate, gamma function, option gamma \\
\(\delta\) & small change, tolerance, Dirac/Kronecker delta \\
\(\varepsilon\) & error, tolerance, strain \\
\(\zeta\) & zeta function, damping ratio \\
\(\eta\) & efficiency, learning rate, parameter \\
\(\theta\) & angle, statistical parameter \\
\(\iota\) & inclusion map \\
\(\kappa\) & curvature, condition number, conductivity \\
\(\lambda\) & eigenvalue, multiplier, rate, wavelength \\
\(\mu\) & mean, measure, viscosity \\
\(\nu\) & frequency, degrees of freedom, kinematic viscosity \\
\(\xi\) & random variable, coordinate \\
\(\pi\) & circle constant \\
\(\rho\) & density, correlation, spectral radius \\
\(\sigma\) & standard deviation, stress, volatility \\
\(\tau\) & time constant, delay, stopping time \\
\(\upsilon\) & specialized parameter or variable \\
\(\phi,\varphi\) & angle, map, potential, function \\
\(\chi\) & chi-square, characteristic quantity \\
\(\psi\) & wave function, stream function \\
\(\omega\) & angular frequency, vorticity \\
\(\Gamma\) & gamma function, circulation \\
\(\Delta\) & change, Laplacian, discriminant \\
\(\Theta\) & asymptotic notation, option theta \\
\(\Lambda\) & diagonal eigenvalue matrix, parameter set \\
\(\Pi\) & product, projection-related notation \\
\(\Sigma\) & summation, covariance matrix \\
\(\Phi\) & normal CDF, function or map \\
\(\Psi\) & wave/state function \\
\(\Omega\) & sample space, domain
\end{tabular}
\end{center}

%%%%%%%%%%%%%%%%%%%%%%%%%%%%%%%%%%%%%%%%%%%%%%%%%%%%%%%%%%%%
\subsection{Conceptual Summary}
\label{subsec:greek-alphabet-summary}
%%%%%%%%%%%%%%%%%%%%%%%%%%%%%%%%%%%%%%%%%%%%%%%%%%%%%%%%%%%%

Greek letters provide a large and flexible collection of symbols used across
mathematics.

The complete alphabet contains twenty-four letters.

Frequently encountered lowercase symbols include

\[
\boxed{
\alpha,\beta,\gamma,\delta,\varepsilon,\theta,\lambda,\mu,\sigma,\phi,\omega.
}
\]

Frequently encountered uppercase symbols include

\[
\boxed{
\Gamma,\Delta,\Theta,\Lambda,\Pi,\Sigma,\Phi,\Psi,\Omega.
}
\]

Certain letters have alternate forms such as

\[
\boxed{
\epsilon,\varepsilon,
\qquad
\phi,\varphi,
\qquad
\theta,\vartheta.
}
\]

The central rule is:

\[
\boxed{
\text{a Greek letter has no universal mathematical meaning by itself}.
}
\]

Its meaning must be defined by the mathematical context.

%%%%%%%%%%%%%%%%%%%%%%%%%%%%%%%%%%%%%%%%%%%%%%%%%%%%%%%%%%%%
\subsection{Section Connections}
\label{subsec:greek-alphabet-connections}
%%%%%%%%%%%%%%%%%%%%%%%%%%%%%%%%%%%%%%%%%%%%%%%%%%%%%%%%%%%%

\chapterconnection{
The Greek Alphabet reference supplements
Section~\ref{sec:mathematical-symbols-notation} by collecting the Greek
symbols used throughout the textbook and identifying many of their common
roles. Greek letters occur throughout calculus, algebra, probability,
statistics, differential equations, optimization, geometry, physics,
continuum mechanics, fluid dynamics, control theory, finance, and economics.
The next section, Number Systems, consolidates the major numerical sets and
their relationships, including natural numbers, integers, rational numbers,
real numbers, complex numbers, and several important extensions and
substructures.
}

%%%%%%%%%%%%%%%%%%%%%%%%%%%%%%%%%%%%%%%%%%%%%%%%%%%%%%%%%%%%
% SECTION SOURCE TRACKING
%%%%%%%%%%%%%%%%%%%%%%%%%%%%%%%%%%%%%%%%%%%%%%%%%%%%%%%%%%%%

%------------------------------------------------------------
% 14.2 Greek Alphabet
%------------------------------------------------------------
%
% Source basis:
%   - Standard Greek alphabet.
%   - Standard mathematical notation conventions.
%   - Consolidated Greek-letter usage from across the textbook.
%
% Used for:
%   - complete lowercase Greek alphabet;
%   - complete uppercase Greek alphabet;
%   - alternate Greek forms;
%   - LaTeX commands;
%   - angles;
%   - parameters;
%   - eigenvalues;
%   - probability notation;
%   - statistical notation;
%   - calculus and analysis notation;
%   - differential-equation notation;
%   - optimization notation;
%   - topology and algebra notation;
%   - number-theory notation;
%   - physics notation;
%   - continuum-mechanics notation;
%   - fluid-dynamics notation;
%   - control-theory notation;
%   - finance notation;
%   - economics notation;
%   - random-variable notation;
%   - tensor indices;
%   - measures;
%   - maps;
%   - common Greek-symbol ambiguities;
%   - symbol-choice conventions.
%
% Notes:
%   - Number Systems is developed next in Section 14.3.
%   - Mathematical Constants follows in Section 14.4.
%
%%%%%%%%%%%%%%%%%%%%%%%%%%%%%%%%%%%%%%%%%%%%%%%%%%%%%%%%%%%%